% !TeX root = ../main.tex

\chapter{Chaos Engineering e Test di Fallimento}
\label{cap:chaos_engineering}

La suite \texttt{ops/chaos\_test\_suite.py} inietta guasti controllati sui componenti critici (database, broker, capacità di flotta) per misurare la reazione del sistema in termini di RTO e capacità di auto-rimediazione. Ogni esperimento segue il ciclo: ipotesi (SLO atteso), blast radius confinato al Deployment in namespace \texttt{default}, iniezione via API K8s o generatore di carico, osservazione su \texttt{chaos\_test\_results.log}, abort condition a 60\,s.

\section{Esperimenti e Risultati}
\label{sec:esperimenti_chaos}

\textbf{(1) Primary Database Outage.} Scaling di \texttt{influxdb} a 0 repliche, attesa 10\,s, ripristino a 1; cronometraggio del ritorno a \texttt{Ready}. \textit{Risultato:} RTO = 12{,}05\,s. Il \texttt{central\_server} ha continuato a servire la dashboard dallo stato in-memory; solo le scritture storiche sono andate perse, coerentemente col pattern di dual-write (§\ref{sec:graceful_degradation}).

\textbf{(2) MQTT Broker Crash.} Deletion del Pod \texttt{mosquitto} via \texttt{delete\_namespaced\_pod}, attesa del nuovo Pod \texttt{Running}. \textit{Risultato:} downtime di scheduling K8s = 0{,}04\,s. La riconnessione dei droni è gestita nativamente da \texttt{paho-mqtt} senza logica custom; la telemetria \textit{at-most-once} (2\,s) tollera le perdite.

\textbf{(3) Load Spike e Scaling Agentico.} Iniezione di 50 ordini con \texttt{client\_simulator.py --stress} (durata $\approx$56\,s), attesa ulteriore di 30\,s per il ciclo decisionale e rilettura delle repliche. \textit{Risultato:} repliche prima = 3, dopo = 6. L'orchestratore agentico ha rilevato lo spike, invocato il tool MCP di scaling e raddoppiato la flotta entro la finestra osservativa ($\approx$86\,s end-to-end), confermando l'efficacia dello scaling agentico per shock di carico controllati.

\begin{table}[ht]
\centering
\begin{tabular}{lll}
\toprule
\textbf{Esperimento} & \textbf{Metrica} & \textbf{Valore} \\
\midrule
InfluxDB outage      & RTO          & 12{,}05\,s \\
MQTT broker crash    & Downtime K8s & 0{,}04\,s \\
Load spike (50 ord.) & $\Delta$ repliche & 3 $\rightarrow$ 6 ($\approx$86\,s) \\
\bottomrule
\end{tabular}
\caption{Run \texttt{chaos\_test\_suite.py} del 23/05/2026.}
\label{tab:chaos_results}
\end{table}

\section{Lezioni Apprese}
\label{sec:lezioni_chaos}

Gli esperimenti confermano la robustezza dei livelli infrastrutturali (Kubernetes + \texttt{paho-mqtt} + dual-write) e validano lo scaling agentico su orizzonti di $\approx$1\,min. Restano due limiti: (i) la latenza del ciclo decisionale ($\approx$30\,s) è ancora elevata per spike sub-minuto, dove un HPA nativo offrirebbe un \emph{fast-path} complementare; (ii) la suite non misura ancora la \textit{loss rate} di telemetria durante l'outage di InfluxDB né la latenza end-to-end sotto stress. Tali estensioni sono discusse nel Capitolo~\ref{cap:conclusioni}.
